% TikZ 图：回旋加速器模型
% 图示正电荷、磁场垂直纸面向里
\begin{tikzpicture}[scale=0.88,>=Latex,line cap=round,line join=round]
  \path[use as bounding box] (-2.75,-3.25) rectangle (2.75,3.10);

  % 两个 D 形盒与窄缝
  \fill[gray!10] (-0.16,2.20) -- (-0.16,-2.20)
    arc (-90:-270:2.20) -- cycle;
  \fill[gray!10] (0.16,2.20) -- (0.16,-2.20)
    arc (-90:90:2.20) -- cycle;
  \draw[thick] (-0.16,2.20) -- (-0.16,-2.20)
    arc (-90:-270:2.20);
  \draw[thick] (0.16,2.20) -- (0.16,-2.20)
    arc (-90:90:2.20);
  \node[font=\small] at (-1.35,2.52) {D 形盒};
  \node[font=\small] at (1.35,2.52) {D 形盒};

  % 磁场：垂直纸面向里
  \foreach \x in {-1.65,-1.05,-0.52,0.52,1.05,1.65} {
    \foreach \y in {-1.55,-0.78,0.78,1.55} {
      \draw[blue!65!black,thick] (\x-0.06,\y-0.06) -- (\x+0.06,\y+0.06);
      \draw[blue!65!black,thick] (\x-0.06,\y+0.06) -- (\x+0.06,\y-0.06);
    }
  }
  \node[blue!65!black,font=\small] at (0,-2.64)
    {$B$ 垂直纸面向里};

  % 一条连续的、半径逐次增大的轨迹
  \draw[orange!90!black,very thick,->] (0,0) arc (-90:90:0.35);
  \draw[orange!90!black,very thick,->] (0,0.70) arc (90:270:0.65);
  \draw[orange!90!black,very thick,->] (0,-0.60) arc (-90:90:0.95);
  \draw[orange!90!black,very thick,->] (0,1.30) arc (90:270:1.25);
  \draw[orange!90!black,very thick,->] (0,-1.20) arc (-90:66:1.55);
  \fill[red] (0,0) circle (2.3pt);

  % 只画一个电场方向示意，避免把两个时刻画在同一图中
  \draw[->,red,very thick] (-0.48,1.82) -- (0.48,1.82);
  \node[red,right,font=\small] at (0.50,1.82) {缝隙电场 $E$};
  \node[font=\small] at (0,2.86)
    {交变电场每隔 $T/2$ 反向一次};
  \node[font=\small] at (0,-3.02)
    {盒内磁场负责转弯，窄缝电场负责做功};
\end{tikzpicture}
